\chapter*{Abstract}
\addcontentsline{toc}{chapter}{Abstract}
Traditional 3D geomodelling techniques face a significant compromise between computational efficiency and geological precision. While traditional methods struggle to simultaneously achieve high visual realism and adapt to hard well data, \gls{dl} algorithms offer rapid sampling capabilities well-suited for large-scale uncertainty quantification. This research assesses the applicability of a \gls{gan}, specifically a modified \gls{fluvgan} architecture, to generate 3D realisations of the real-world \gls{dsst} geothermal reservoir while conditioned to sparse real-world data. Robust synthetic training datasets were curated using the process-based stochastic modelling software FLUMY. 

The results demonstrate that while a properly tuned 3-class \gls{gan} model successfully captures macro-scale 3D spatial dependencies and replicates dominant volumetric proportions, its performance is heavily bound by hyperparameter configurations. Stable adversarial dynamics were achieved by expanding network depth via double residual blocks and skip connections, combined with targeted R1 gradient regularization to preserve high-frequency spatial gradients. Expanding to a 9-class high-fidelity model or a model trained on a multi-prior dataset triggered severe optimization bottlenecks and ensemble mode collapse due to extreme categorical class imbalances within the training data. 

To bridge this to practical field application, a two-step inversion workflow combining latent space optimization with pivotal tuning was deployed to anchor the synthetic realisations to hard physical constraints from the DEL-GT-01 and DEL-GT-02-S2 vertical well trajectories. This approach successfully matched the asymmetric local observations, yielding localized accuracy scores ranging from 0.78 to 0.92. However, true geological plausibility remains bottlenecked by topological artifacts—such as broken mudplugs and isolated voxels—in statistically suppressed minority facies, alongside a deflated spatial variance across the realization ensemble. This study concludes that while \glspl{gan} offer an efficient pathway for ensemble-based uncertainty quantification, successful full-field subsurface deployment depends on advanced loss formulations and architectural upscaling to resolve fine-scale topological features without sacrificing stochastic diversity.